\begin{abcliste}
\abc With $\ssc{E}{kin} = \frac{1}{2}m_n v^2$ (the energy in joules):
\begin{align}
 v &= \vF \\
   &= \sqrt{\frac{2\times\EkX\times\nce}{\ncmn}} \\
   &= \v \approx \result{\vP{3}{2}}
\end{align}

\abc
\begin{align}
 \lambda &= \frac{h}{m_n v} = \laF \\
   &= \frac{\nch}{\sqrt{2\times\ncmn\times\EkX\times\nce}} \\
   &= \la \approx \result{\laP{-9}{2}}
\end{align}
about the spacing of the atoms, just right for diffraction.

\abc
\begin{align}
 \theta &= \thtF = \arcsin\frac{\la}{2\times\d} = \tht \approx \result{\thtP{0}{2}}
\end{align}
(Forgetting the factor 2 in Bragg's condition is the typical mistake.) Neutrons are uncharged and see the nuclei and their magnetism: they show things X-rays miss, for example where the hydrogen atoms are.
\end{abcliste}
